\begin{table}[htbp]
\centering
\small
\caption{V3 独立ファミリー再現性マトリックス：探索段階（Qwen）で得られた知見の独立3ファミリー（Llama 3.2, Gemma 3, OLMo 2）における検証結果総括。各仮説はValenceおよびArousalの双方が基準を満たした場合のみ支持（$\checkmark$）とされる。}
\label{tab:v3_confirmatory_matrix}
\begin{tabular}{l cccc c}
\toprule
\textbf{Family} & \textbf{H1 (Dissociation)} & \textbf{H2 (Sufficiency)} & \textbf{H3 (Mediation)} & \textbf{H4 (Temporal Contrast)} & \textbf{All Confirmed?} \\
\midrule
Llama 3.2      & $\times$     & $\times$     & $\times$     & \checkmark   & \textbf{NO} \\
Gemma 3        & $\times$     & $\times$     & $\times$     & \checkmark   & \textbf{NO} \\
OLMo 2         & $\times$     & $\times$     & $\times$     & \checkmark   & \textbf{NO} \\
\bottomrule
\end{tabular}
\vspace{1ex}
\begin{minipage}{\linewidth}
\footnotesize
\textbf{Note:} 全4仮説の総合判定マトリックス。各セルは該当仮説においてValenceおよびArousalの双方が事前登録基準（95\% CI）を満たした場合に合致（$\checkmark$）とする。全モデルファミリーにおいてH1--H3は棄却され、H4のみが支持されたため、全体結論として情動表現の機能的・因果的活用仮説は支持されなかった（All Confirmed = NO）。
\end{minipage}
\end{table}